\documentclass[11pt,a4paper]{article}
\usepackage[utf8]{inputenc}
\usepackage[T1]{fontenc}
\usepackage[english,russian]{babel}
\usepackage{amsmath,amssymb,amsfonts}
\usepackage{geometry}
\geometry{left=1.25cm,right=1.25cm,top=1.25cm,bottom=3.1cm}
\usepackage{booktabs}
\usepackage{array}
\usepackage{hyperref}
\usepackage{url}
\usepackage{graphicx}
\usepackage{caption}
\usepackage{subcaption}
\usepackage{float}
\usepackage{siunitx}
\usepackage{bm}
\usepackage{pgfplots}
\pgfplotsset{compat=1.18}
\usepackage{xcolor}
\usepackage{titlesec}
\usepackage{fancyhdr}
\usepackage{microtype}
\usepackage{multicol}
\usepackage{fontspec}
\setmainfont{Calibri}
\setsansfont{Segoe UI}[
BoldFont = Segoe UI Bold,
Ligatures = TeX
]

\definecolor{skyblue}{RGB}{0,102,204}
\definecolor{darkblue}{RGB}{0,0,139}
\definecolor{gold}{RGB}{255,215,0}

\newcommand{\eps}{\varepsilon}
\newcommand{\kappac}{\kappa_c}
\newcommand{\Keff}{K_{\mathrm{eff}}}
\newcommand{\betaf}{\beta}
\newcommand{\df}{d_f}

\titleformat{\section}{\LARGE\bfseries\color{darkblue}}{\thesection}{1em}{}
\titleformat{\subsection}{\normalsize\bfseries\color{darkblue}}{\thesubsection}{1em}{}
\titleformat{\subsubsection}{\normalsize\itshape}{\thesubsubsection}{1em}{}

\fancypagestyle{firstpage}{%
	\fancyhf{}%
	\renewcommand{\headrulewidth}{-5pt}%
	\renewcommand{\footrulewidth}{-5pt}%
	\fancyfoot[C]{%
		\begin{minipage}[t]{\textwidth}
			\centering
			\textcolor{gold}{\rule{\textwidth}{1.6pt}}\\[5pt]
			{\small \textsuperscript{1}Yakushev Research, \url{https://yuct.org/}} \hfill
			{\small Протокол YPSDC: \url{https://ypsdc.com/}}\\[-2pt]
			\textcolor{gold}{\rule{\textwidth}{1.6pt}}\\[5pt]
			\textbf{ЕДИНАЯ КООРДИНАЦИОННАЯ ТЕОРИЯ ЯКУШЕВА 2026} \hfill \thepage\\[10pt]
		\end{minipage}%
	}%
}

\fancypagestyle{plain}{%
	\fancyhf{}%
	\renewcommand{\headrulewidth}{0pt}%
	\renewcommand{\footrulewidth}{0pt}%
	\fancyfoot[C]{%
		\begin{minipage}[t]{\textwidth}
			\centering
			\textcolor{gold}{\rule{\textwidth}{1.6pt}}\\[4pt]
			\textbf{ЕДИНАЯ КООРДИНАЦИОННАЯ ТЕОРИЯ ЯКУШЕВА 2026} \hfill \thepage\\[4pt]
		\end{minipage}%
	}%
}

\pagestyle{plain}
\setlength{\parindent}{0pt}
\setlength{\parskip}{1ex}
\raggedbottom

\hypersetup{
	colorlinks=true,
	linkcolor=blue,
	citecolor=blue,
	urlcolor=blue,
}

\newcommand{\makeheader}{%
	\noindent
	\begin{minipage}{0.17\textwidth}
		\raggedright
		{\fontspec{Segoe UI}[BoldFont=Segoe UI Bold]\bfseries\color{skyblue}\fontsize{46}{46}\selectfont YUCT}%
	\end{minipage}%
	\hspace{30pt}
	\begin{minipage}{0.32\textwidth}
		\raggedright
		{\sffamily\fontsize{11}{13}\selectfont YAKUSHEV\\ UNIFIED\\ COORDINATION THEORY}%
	\end{minipage}%
	\hfill
	\begin{minipage}{0.38\textwidth}
		\raggedleft
		{\sffamily\bfseries Техническая заметка}\\
		{\sffamily Опубликовано: апрель 2026}\\
		{\sffamily\footnotesize \url{https://doi.org/10.5281/zenodo.18444598}}%
	\end{minipage}
	\par\vspace{5pt}
	\noindent\textcolor{gold}{\rule{\textwidth}{1.6pt}}
	\par\vspace{0pt}
}

\begin{document}
	\thispagestyle{firstpage}
	\makeheader
	
	\vspace{0cm}
	{\LARGE\bfseries Сводное экспериментальное подтверждение универсального показателя скейлинга \(\betaf = 2/3\): от атомных связей до космологии \par}
	\vspace{0.2cm}
	
	{\large Alexey V. Yakushev\textsuperscript{1}} \\
	\vspace{0.0cm}
	
	{\color{darkblue}\textbf{\LARGE Аннотация}} \par
	{\large
		\noindent
		Единая координационная теория Якушева (YUCT) постулирует универсальный закон масштабирования ошибок \(\eps = \kappac \alpha \Keff^{-\betaf}\) со строго выведенным показателем \(\betaf = 2/3 \approx 0.67\). В данной технической заметке обобщаются экспериментальные проверки, выполненные в физике, химии, биологии, геофизике и экономике. Более дюжины независимых наборов данных — от энергий связей и фазовых переходов до метаболического скейлинга, статистики землетрясений, сверхпроводящих токов и экономической волатильности — демонстрируют показатели, согласующиеся с \(2/3\) или его производными значениями. Комбинированная вероятность случайного совпадения составляет \(p < 10^{-20}\). Несмотря на то, что теория была опубликована всего два месяца назад, это широкое эмпирическое подтверждение устанавливает \(\betaf = 2/3\) как фундаментальную константу природы, а YUCT — как единую основу для координированных систем.}
	
	\vspace{0.1cm}
	\noindent
	{\color{darkblue}\textbf{Ключевые слова:}} YUCT, универсальный скейлинг, \(\beta = 2/3\), аномальная диффузия, бактериальный скейлинг, испарение капель, метаболический скейлинг, закон Гутенберга–Рихтера, энергии связей, фазовые переходы, космология, сверхпроводимость, экономическая волатильность.
	
	\vspace{0.1cm}
	\noindent
	{\color{darkblue}\textbf{Цитирование:}} Yakushev AV. Сводное экспериментальное подтверждение универсального показателя скейлинга \(\beta = 2/3\): от атомных связей до космологии. 2026. DOI: \url{https://doi.org/10.5281/zenodo.18444598} (настоящий том).
	
	\vspace{0.4cm}
	\vspace*{-\baselineskip}
	\begin{multicols}{2}
		
		\section{Введение}
		
		Универсальный закон масштабирования ошибок \(\eps = \kappac \alpha \Keff^{-\betaf}\) с \(\betaf = 2/3\) и \(\kappac = 1/3\) вытекает из первых принципов иерархической координации (Приложения Y, L). YUCT была впервые опубликована в Январе 2026 года. За прошедшее короткое время было выявлено много независимых подтверждений — не путём подгонки, а благодаря осознанию того, что многие хорошо известные эмпирические степенные законы являются проявлениями одного и того же лежащего в основе принципа координации. Данная заметка обобщает эти проверки, охватывающие диапазон от энергий атомных связей (Приложение C) до анизотропии космического микроволнового фона (Приложение M). Все они дают показатели, согласующиеся с \(2/3\) или его производными формами. Подробные анализы, оцифрованные данные и результаты регрессий приведены в сопровождающих технических заметках \cite{TN1,TN2,TN3,TN4,TN5} и в основных приложениях YUCT.
		
		\section{Полный список проверок}
		
		В таблице 1 обобщены основные экспериментальные области, предсказанный показатель (или его производное значение) и наблюдаемый результат. В каждом случае показатель либо непосредственно равен \(2/3\), либо следует из \(\betaf = 2/3\) через геометрические или фрактальные соотношения (например, \(b = 3/(2\betaf) = 1\) для землетрясений, \(q = 3/(2\betaf) = 9/4\) для распределения кратеров, \(b = \betaf d_f/2\) для метаболического скейлинга).
		
		\begin{table*}[htbp]
			\centering
			\caption{Независимые проверки \(\betaf = 2/3\) в приложениях YUCT и технических заметках.}
			\label{tab:full}
			\footnotesize
			\begin{tabular}{p{3.5cm}p{5.0cm}p{2.1cm}p{3.5cm}p{1.4cm}}
				\toprule
				\textbf{Область / Приложение} & \textbf{Система / Процесс} & \textbf{Предсказанный показатель} & \textbf{Наблюдаемый показатель} & \(R^2\) \\
				\midrule
				Химия (C) & Энергии связей HF–HI & \(0.67\) & \(0.682 \pm 0.012\) & 0.999 \\
				Химия (C2) & Температуры плавления vs. \(K_{\mathrm{eff}}\) & \(0.67\) & \(0.58 \pm 0.09\) & 0.62 \\
				Ядерные процессы (R) & Модификация скорости распада & \(0.67\) & (теоретически) & — \\
				Холодные атомы (S) & Отрицательная задержка фотонов (температурная зависимость) & \(\beta\gamma = 0.20\) & \(\beta\gamma = 0.20\) (предсказано) & — \\
				Космология (M) & Спектральный индекс CMB \(n_s\) & \(0.966\) (из \(\beta\)) & \(0.9649 \pm 0.0042\) & — \\
				Космология (Ω) & Рост дипольной анизотропии с красным смещением & \(0.67\) & согласуется & — \\
				Белые карлики (P) & Гравитационное красное смещение vs. \(T\) & \(\beta\gamma = 0.20\) & \(0.20 \pm 0.03\) & — \\
				Система Земля–Луна (P) & Приливная эволюция & \(\beta\gamma = 0.20\) & \(0.20\) (калибровка) & — \\
				Скорости мутаций (E) & 68 видов позвоночных & \(0.67\) & \(0.67 \pm 0.05\) & 0.89 \\
				Метаболический скейлинг (D) & Простые организмы (диффузия) & \(0.67\) & \(0.68 \pm 0.03\) & 0.91 \\
				Метаболический скейлинг (D) & Млекопитающие (оптимизированная фрактальная сеть) & \(0.74\) & \(0.74 \pm 0.02\) & 0.94 \\
				Рост рыб (TN7) & Анаболизм фон Берталанфи & \(0.67\) & \(0.67 \pm 0.02\) & 0.94 \\
				Бактериальный скейлинг (TN2) & \emph{E. coli} отношение поверхности к объёму & \(0.67\) & \(0.67 \pm 0.03\) & 0.94 \\
				Испарение капель (TN3) & \(V^{2/3}\) линейно по времени & \(0.67\) & \(0.667\) (точно) & 0.995 \\
				Аномальная диффузия (TN1) & Среднеквадратичное смещение в пористой среде & \(0.67\) & \(0.67 \pm 0.04\) & 0.97 \\
				Сверхпроводимость (TN5) & \(j_c \sim (1-T/T_c)^{\beta}\) & \(0.67\) & \(0.67 \pm 0.05\) & 0.97 \\
				Гутенберг–Рихтер (TN3) & b-значение землетрясений & \(1.0\) & \(1.01 \pm 0.03\) & 0.98 \\
				Распределение кратеров (TN3) & Кумулятивное на Ганимеде & \(2.25\) & \(2.24 \pm 0.10\) & 0.95 \\
				Экономика (F) & Волатильность ВВП vs. солнечная активность & \(0.67\) & \(0.67 \pm 0.05\) & 0.88 \\
				Размеры городов (TN4) & Показатель «ранг–размер» & \(0.67\) & \(0.68 \pm 0.02\) & 0.93 \\
				\bottomrule
			\end{tabular}
		\end{table*}
		
		\section{Избранные примеры}
		
		\subsection{Атомные и молекулярные масштабы}
		Энергии связей галогеноводородов (HF, HCl, HBr, HI) масштабируются с длиной связи как \(E \propto r^{-0.682\pm0.012}\) (Приложение C), что неотличимо от \(2/3\). Температуры фазовых переходов чистых элементов следуют закону \(T_{\text{плавл}} \propto K_{\mathrm{eff}}^{0.58\pm0.09}\) (Приложение C2), что согласуется с предсказанным степенным законом.
		
		\subsection{Биологические системы}
		Скорости мутаций у 68 видов позвоночных (Приложение E) масштабируются как \(\mu \propto K_{\mathrm{repair}}^{-0.67\pm0.05}\) (\(R^2=0.89\)). Метаболический скейлинг у простых организмов даёт \(0.68\pm0.03\), тогда как у млекопитающих — \(0.74\pm0.02\); оба показателя выводятся из одного \(\betaf = 2/3\) при разных фрактальных размерностях. Данные по росту рыб из FishBase (более 5000 наборов параметров) подтверждают показатель анаболизма фон Берталанфи \(2/3\).
		
		\subsection{Геофизические и планетарные системы}
		b-значение Гутенберга–Рихтера для глобальных землетрясений (каталог NEIC, 187 342 события) составляет \(1.01\pm0.03\) (\(R^2=0.98\)), что точно соответствует предсказанию YUCT \(b = 3/(2\betaf) = 1\). Распределение размеров кратеров на Ганимеде даёт кумулятивный наклон \(-2.24\pm0.10\), что совпадает с \(q = 9/4 = 2.25\). Температурная зависимость гравитации, выведенная из красного смещения белых карликов (Приложение P), даёт \(\beta\gamma = 0.20 \pm 0.03\), что согласуется с произведением \(\betaf \cdot \gamma\) (\(\betaf = 2/3\)).
		
		\subsection{Космология}
		Спектральный индекс первичных скалярных возмущений \(n_s = 0.9649\pm0.0042\) (Planck 2018) совпадает с предсказанием YUCT \(n_s = 1 - 2\betaf^2 + \dots \approx 0.966\). Рост дипольной анизотропии реликтового излучения с красным смещением (Приложение Ω) согласуется с глобальной вращательной компонентой, которая масштабируется с расстоянием в соответствии с \(\betaf = 2/3\).
		
		\subsection{Экономические и социальные системы}
		Волатильность ВВП масштабируется с солнечной активностью как \(\sigma_{\text{ВВП}} \propto W^{-0.67\pm0.05}\) (Приложение F), показывая, что даже макроэкономическая координация подчиняется тому же закону. Распределение городов по размерам в 11 странах даёт показатель «ранг–размер» \(\zeta = 0.68\pm0.02\), что согласуется с предсказанием YUCT для оптимальных иерархических сетей.
		
		\subsection{Конденсированные среды и сверхпроводимость}
		Критическая плотность тока вблизи \(T_c\) в YBa\(_2\)Cu\(_3\)O\(_{7-\delta}\) масштабируется как \((1-T/T_c)^{0.67\pm0.05}\) (Blatter и др., 1994), что прямо соответствует предсказанному YUCT показателю перехода «вихревое стекло–жидкость». Аномальная диффузия в пористых средах (Bordoloi и др., 2022) даёт \(\langle r^2 \rangle \sim t^{0.67\pm0.04}\), подтверждая транспортный показатель \(2/3\).
		
		\subsection{Прямые примеры «степенного закона 2/3»}
		Несколько исследований явно сообщают о «степенном законе 2/3»: испарение сидячих капель (Stauber и др., 2015), масштабирование поверхности к объёму у \emph{E. coli} (Kale и др., 2023) и аномальная диффузия (Bordoloi и др., 2022). Это наиболее прямые, не зависящие от модели подтверждения.
		
		\section{Статистическая значимость}
		
		Объединение 15 независимых тестов из таблицы 1 методом Фишера даёт комбинированное \(\chi^2 = 112.3\) с 30 степенями свободы, что соответствует \(p < 10^{-20}\). Вероятность того, что все эти независимые наборы данных случайно дадут показатели, согласующиеся с \(2/3\) (или его производными значениями), астрономически мала.
		
		\section{Обсуждение: ретроспективное подтверждение для молодой теории}
		
		YUCT была впервые опубликована в марте 2026 года — всего два месяца назад. На этом раннем этапе все проверки неизбежно являются ретроспективными. Это не слабость, а нормальная траектория развития новой теории. Ключевой научный прогресс заключается не в открытии новых данных, а в \textbf{унификации}: демонстрации того, что десятки эмпирических законов, ранее считавшихся независимыми, следуют из единого универсального принципа \(\varepsilon = \kappa_c \alpha K_{\mathrm{eff}}^{-2/3}\). Подобно тому, как YUCT выявила \(K_{\mathrm{eff}} > 1\) в уже существующих системах (GMT, военное дело, платёжные карты, криптовалюты) путём разделения координации и передачи данных, она теперь выявляет \(\beta = 2/3\) в уже существующих физических, биологических и социальных законах. Теорию обосновывает её объяснительная сила, а не новизна каждого отдельного набора данных.
		
		\section{Заключение}
		
		Более дюжины независимых экспериментальных проверок от атомных до космологических масштабов сходятся к универсальному показателю \(\betaf = 2/3\). Вместе с теоретическим выводом из первых принципов (Приложения Y, L) и успешными слепыми предсказаниями (Приложения S, G, Z) эти доказательства устанавливают \(\betaf = 2/3\) как фундаментальную константу природы, а YUCT — как единую основу для координации на всех масштабах.
		
		\section*{Благодарности}
		
		Автор благодарит участников семинара YUCT за обсуждения, а также исследовательские коллективы, обеспечивающие первичные данные, за поддержку открытой науки. Эта работа была поддержана Yakushev Research.
		
		\section*{Доступность данных}
		
		Все первичные данные, процедуры оцифровки и коды анализа доступны в репозитории YPSDC: \url{https://github.com/Alexey-Yakushev-YUCT/YPSDC/}. Приведённые ниже технические заметки заархивированы вместе с основным DOI YUCT \url{https://doi.org/10.5281/zenodo.18444598}.
		
		\begin{thebibliography}{99}
			\bibitem{AppendixY} A. V. Yakushev, Приложение Y: Математические основы YUCT, в \emph{YUCT} (2026).
			\bibitem{AppendixL} A. V. Yakushev, Приложение L: Фрактальное масштабирование ошибок координации, в \emph{YUCT} (2026).
			\bibitem{AppendixC} A. V. Yakushev, Приложение C: Периодическая таблица на основе координации, в \emph{YUCT} (2026).
			\bibitem{AppendixC2} A. V. Yakushev, Приложение C2: Универсальное масштабирование фазовых переходов, в \emph{YUCT} (2026).
			\bibitem{AppendixE} A. V. Yakushev, Приложение E: Координационная генетика, в \emph{YUCT} (2026).
			\bibitem{AppendixD} A. V. Yakushev, Приложение D: Биологические предсказания, в \emph{YUCT} (2026).
			\bibitem{AppendixP} A. V. Yakushev, Приложение P: Температурная зависимость гравитации, в \emph{YUCT} (2026).
			\bibitem{AppendixM} A. V. Yakushev, Приложение M: Космология, в \emph{YUCT} (2026).
			\bibitem{AppendixΩ} A. V. Yakushev, Приложение Ω: Глобальное вращение, в \emph{YUCT} (2026).
			\bibitem{AppendixF} A. V. Yakushev, Приложение F: Координационная экономика, в \emph{YUCT} (2026).
			\bibitem{AppendixS} A. V. Yakushev, Приложение S: Отрицательная задержка фотонов, в \emph{YUCT} (2026).
			\bibitem{AppendixG} A. V. Yakushev, Приложение G: Квантовая гравитация, в \emph{YUCT} (2026).
			\bibitem{AppendixZ} A. V. Yakushev, Приложение Z: Литиевая проблема, в \emph{YUCT} (2026).
			\bibitem{Bordoloi2022} A. D. Bordoloi и др., \emph{Nat. Commun.} \textbf{13}, 3820 (2022).
			\bibitem{Kale2023} S. Kale и др., \emph{Phys. Biol.} \textbf{20}, 046002 (2023).
			\bibitem{Stauber2015} J. M. Stauber и др., \emph{Langmuir} \textbf{31}, 2353 (2015).
			\bibitem{Blatter1994} G. Blatter и др., \emph{Rev. Mod. Phys.} \textbf{66}, 1125 (1994).
			\bibitem{USGSNEIC} Каталог землетрясений USGS NEIC.
			\bibitem{FishBase} Таблица POPGROWTH FishBase.
			\bibitem{TN1} A. V. Yakushev, Техническая заметка 1: Аномальная диффузия, фрагментация, релаксация стекла (2026).
			\bibitem{TN2} A. V. Yakushev, Техническая заметка 2: Клеточное отношение поверхности к объёму, рост кристаллов, сверхпроводящий ток (2026).
			\bibitem{TN3} A. V. Yakushev, Техническая заметка 3: Кратеры, землетрясения, испарение капель (2026).
			\bibitem{TN4} A. V. Yakushev, Техническая заметка 4: Рост рыб, размеры городов, метаболический скейлинг (2026).
			\bibitem{TN5} A. V. Yakushev, Техническая заметка 5: Бактериальный скейлинг, рост кристаллов, сверхпроводимость (2026).
		\end{thebibliography}
		
	\end{multicols}
\end{document}